\chapter{Holomorphic Vector Bundles}\label{holomorphic-bundles}
\chapterstatus{complete}

\section{Introduction}\label{holomorphic-bundles:section-introduction}

Holomorphic line bundles are the objects that encode divisors, projective embeddings and, through their first Chern classes, the
$(1,1)$-part of cohomology. The exponential sequence (Theorem~\ref{holomorphic-bundles:theorem-exponential}) computes the Picard
group in terms of cohomology and already contains the case $p = 1$ of the Hodge conjecture: the classes of divisors are exactly the
integral classes of type $(1,1)$ (Chapter~\ref{lefschetz-11}). References: \cite[Chapter 1]{griffiths-harris},
\cite[Chapters 2--4]{huybrechts-complex}, \cite[Chapter 4]{voisin-hodge-1}.

\noindent\textbf{Prerequisites.} Chapter~\ref{complex-manifolds} (Complex Manifolds), Chapter~\ref{sheaf-cohomology} (Sheaf
Cohomology) and Chapter~\ref{characteristic-classes} (Characteristic Classes).

$X$ denotes a complex manifold of dimension $n$.

\section{Holomorphic bundles and sheaves}\label{holomorphic-bundles:section-bundles}

\begin{definition}\label{holomorphic-bundles:definition-holomorphic-bundle}
A \emph{holomorphic vector bundle} of rank $r$ on $X$ is a complex manifold $E$ with a holomorphic map $\pi \colon E \to X$ and $\CC$-vector space structures on the fibres, locally biholomorphic to
$U \times \CC^r$ over $U$ and linear on fibres; equivalently, it is given by holomorphic transition functions $g_{ij} \colon U_i \cap U_j \to \GL_r(\CC)$ satisfying the cocycle condition
(Proposition~\ref{vector-bundles:proposition-cocycles}). Its sheaf of holomorphic sections is a locally free $\mathcal{O}_X$-module of rank $r$, and this is an equivalence between holomorphic bundles and
locally free sheaves, as in Proposition~\ref{vector-bundles:proposition-sheaf-of-sections}.
\end{definition}

\begin{example}\label{holomorphic-bundles:example-bundles}
The holomorphic tangent bundle $T^{1,0}X$, with transition functions the holomorphic Jacobians; the bundles $\Omega^p_X = \Lambda^p(T^{1,0}X)^\vee$ of holomorphic $p$-forms and the canonical bundle
$K_X = \Omega^n_X$; the tautological bundle $\mathcal{O}(-1)$ on $\CC\PP^n$ (Exercise~\ref{vector-bundles:exercise-tautological}), its dual $\mathcal{O}(1)$ and the powers
$\mathcal{O}(k) = \mathcal{O}(1)^{\otimes k}$. The holomorphic sections of $\mathcal{O}(k)$ on $\CC\PP^n$ are the homogeneous polynomials of degree $k$ for $k \geq 0$, and $0$ for $k < 0$
(Exercise~\ref{holomorphic-bundles:exercise-sections-ok}).
\end{example}

\begin{proposition}\label{holomorphic-bundles:proposition-picard}
The isomorphism classes of holomorphic line bundles on $X$ form a group $\Pic(X)$ under $\otimes$, and
\[
\Pic(X) \cong H^1(X, \mathcal{O}_X^*),
\]
where $\mathcal{O}_X^*$ is the sheaf of nowhere vanishing holomorphic functions.
\end{proposition}

\begin{proof}
As in Proposition~\ref{vector-bundles:proposition-cocycles} and Theorem~\ref{vector-bundles:theorem-line-bundles}, line bundles correspond to \v{C}ech cocycles with values in $\mathcal{O}_X^*$ modulo
coboundaries, and $\check{H}^1 = H^1$ (Remark~\ref{sheaf-cohomology:remark-cech-limit}); tensor product corresponds to multiplication of cocycles.
\end{proof}

\begin{theorem}[Exponential sequence]\label{holomorphic-bundles:theorem-exponential}
The sequence of sheaves $0 \to \ZZ_X \to \mathcal{O}_X \xrightarrow{\exp(2\pi i\,\cdot)} \mathcal{O}_X^* \to 0$ is exact, and its long exact sequence reads
\[
\cdots \to H^1(X, \ZZ) \to H^1(X, \mathcal{O}_X) \to \Pic(X) \xrightarrow{\ c_1\ } H^2(X, \ZZ) \to H^2(X, \mathcal{O}_X) \to \cdots
\]
The map $c_1$ is the first Chern class of the underlying smooth line bundle (Theorem~\ref{vector-bundles:theorem-line-bundles}). In particular, for compact $X$:
\begin{enumerate}
\item $c_1$ is surjective onto the kernel of $H^2(X, \ZZ) \to H^2(X, \mathcal{O}_X)$;
\item its kernel $\Pic^0(X)$ is $H^1(X, \mathcal{O}_X)/\operatorname{im}H^1(X, \ZZ)$;
\item if $H^1(X, \mathcal{O}_X) = H^2(X, \mathcal{O}_X) = 0$, then $c_1 \colon \Pic(X) \to H^2(X, \ZZ)$ is an isomorphism.
\end{enumerate}
\end{theorem}

\begin{proof}
Exactness: the kernel of $\exp(2\pi i \cdot)$ is the constant sheaf $\ZZ$, and surjectivity holds on stalks because a nonvanishing holomorphic function on a small ball has a holomorphic logarithm (integrate
$df/f$, which is possible on a simply connected domain by Theorem~\ref{complex-analysis:theorem-cauchy-homotopy} applied in one variable at a time, or note that the ball is contractible and use the same
argument as in Theorem~\ref{vector-bundles:theorem-line-bundles}). The long exact sequence is Theorem~\ref{homological-algebra:theorem-long-exact-sequence}, using
$\Pic(X) \cong H^1(X, \mathcal{O}_X^*)$. That the connecting map is the topological first Chern class follows by comparing with the corresponding sequence for smooth functions, into which this one maps, and
using that the sheaf of smooth functions is acyclic. The three consequences are exactness of the long sequence.
\end{proof}

\begin{remark}\label{holomorphic-bundles:remark-picard-variety}
When $X$ is compact K\"ahler, the Hodge decomposition (Chapter~\ref{hodge-decomposition}) shows that $H^1(X, \ZZ)$ maps onto a lattice in the complex vector space $H^1(X, \mathcal{O}_X)$ of dimension
$b_1/2$, so $\Pic^0(X)$ is a compact complex torus, the \emph{Picard variety}, and the image of $c_1$, the \emph{N\'eron--Severi group}, is the group of integral classes of type $(1,1)$
(Chapter~\ref{lefschetz-11}).
\end{remark}

\section{Divisors}\label{holomorphic-bundles:section-divisors}

\begin{definition}\label{holomorphic-bundles:definition-divisor}
A \emph{divisor} on $X$ is a locally finite formal sum $D = \sum_ia_iZ_i$ with $a_i \in \ZZ$ and $Z_i$ irreducible analytic hypersurfaces; they form the group $\mathrm{Div}(X)$. A \emph{meromorphic function}
is a section of the sheaf $\mathcal{M}_X$ of quotients $f/g$ of holomorphic functions with $g \not\equiv 0$; the order $\operatorname{ord}_Z(f) \in \ZZ$ of a meromorphic function along an irreducible
hypersurface $Z$ is defined by writing $f = u\,h^m$ near a regular point of $Z$, where $h$ is a local defining function of $Z$ and $u$ is a unit, using that $\mathcal{O}_{X,x}$ is a unique factorisation domain
(Theorem~\ref{several-complex-variables:theorem-local-ring}). The divisor of $f$ is $(f) = \sum_Z\operatorname{ord}_Z(f)\,Z$, and divisors of this form are \emph{principal}.
\end{definition}

\begin{proposition}\label{holomorphic-bundles:proposition-divisor-line-bundle}
There is a homomorphism $\mathrm{Div}(X) \to \Pic(X)$, $D \mapsto \mathcal{O}(D)$, defined by taking local defining functions $f_i$ for $D$ on a cover and the cocycle $g_{ij} = f_i/f_j \in \mathcal{O}^*$. Its kernel
is the group of principal divisors, and $\mathcal{O}(D)$ has a canonical meromorphic section with divisor $D$; for $D$ effective (all $a_i \geq 0$) this section is holomorphic and vanishes exactly on $D$.
Conversely a nonzero meromorphic section of a line bundle $L$ has a divisor $D$ with $L \cong \mathcal{O}(D)$.
\end{proposition}

\begin{proof}
The functions $f_i/f_j$ are nowhere vanishing and holomorphic on overlaps, because two local defining functions of the same divisor differ by a unit
(Theorem~\ref{several-complex-variables:theorem-local-ring}), and they clearly form a cocycle; changing the $f_i$ changes the cocycle by a coboundary. The local functions $f_i$ glue to a meromorphic section
of $\mathcal{O}(D)$, whose divisor is $D$. If $\mathcal{O}(D)$ is trivial, the cocycle is a coboundary $f_i/f_j = u_i/u_j$, so the $f_i/u_i$ glue to a global meromorphic function with divisor $D$; conversely,
principal divisors give trivial cocycles. The last statement follows by taking the divisor of the section and comparing cocycles.
\end{proof}

\begin{example}\label{holomorphic-bundles:example-divisors-projective}
On $\CC\PP^n$ every irreducible hypersurface is the zero set of an irreducible homogeneous polynomial, and $\mathcal{O}(D) \cong \mathcal{O}(\deg D)$; the map $\deg \colon \mathrm{Div}(\CC\PP^n) \to \ZZ$
induces the isomorphism $\Pic(\CC\PP^n) \cong \ZZ$, which also follows from Theorem~\ref{holomorphic-bundles:theorem-exponential} since $H^1(\CC\PP^n, \mathcal{O}) = H^2(\CC\PP^n, \mathcal{O}) = 0$
(Chapter~\ref{dolbeault}). On a complex torus $X = \CC^n/\Lambda$, by contrast, $H^1(X, \mathcal{O}_X) \neq 0$ and $\Pic^0(X)$ is a nontrivial complex torus, the dual torus.
\end{example}

\section{Hermitian metrics and the Chern connection}\label{holomorphic-bundles:section-chern-connection}

\begin{theorem}\label{holomorphic-bundles:theorem-chern-connection}
Let $E$ be a holomorphic vector bundle with a Hermitian metric $h$. There is a unique connection $\nabla$ on $E$, the \emph{Chern connection}, that is metric and satisfies $\nabla^{0,1} = \bar\partial$. In a
holomorphic frame with $H_{ij} = h(e_i, e_j)$ one has $\theta = \overline{H}^{-1}\partial\overline{H}$ and curvature $F = \bar\partial\theta$, a $(1,1)$-form with values in $\End E$. For a line bundle with local
weight $h = e^{-\varphi}$ in a holomorphic frame, $F = \partial\bar\partial\varphi$ and
\[
c_1(L) = \left[\frac{i}{2\pi}F\right] = \left[\frac{i}{2\pi}\partial\bar\partial\varphi\right] \in H^2(X; \RR),
\]
a class of type $(1,1)$.
\end{theorem}

\begin{proof}
Write $\nabla = \nabla^{1,0} + \nabla^{0,1}$. The condition $\nabla^{0,1} = \bar\partial$ fixes the $(0,1)$-part, and the metric condition
$d\,h(s, t) = h(\nabla s, t) + h(s, \nabla t)$ splits into its $(1,0)$- and $(0,1)$-parts, the first of which reads $\partial h(s, t) = h(\nabla^{1,0}s, t) + h(s, \bar\partial t)$ and determines $\nabla^{1,0}$; this
proves uniqueness, and the resulting formula in a holomorphic frame, $\theta = \overline{H}^{-1}\partial\overline{H}$, defines a connection with the two properties, which proves existence. Then
$F = d\theta + \theta \wedge \theta$ has no $(0,2)$-part because $\bar\partial^2 = 0$ and $\nabla^{0,1} = \bar\partial$, and no $(2,0)$-part by the metric condition; the formula $F = \bar\partial\theta$ follows. For a
line bundle, $H = e^{-\varphi}$ and $\theta = \partial\log H^{-1}\cdot(-1) = -\partial\varphi$ up to sign conventions, so $F = \bar\partial\theta = \partial\bar\partial\varphi$ after simplification, and the Chern class is
computed by Chern--Weil theory (Proposition~\ref{characteristic-classes:proposition-chern-properties}).
\end{proof}

\begin{definition}\label{holomorphic-bundles:definition-positive}
A Hermitian line bundle $(L, h)$ is \emph{positive} if the real $(1,1)$-form $\frac{i}{2\pi}F$ is positive definite, that is, if it is the fundamental form of a Hermitian metric on $X$
(Chapter~\ref{kahler}); $L$ is positive if it admits such a metric. A line bundle is \emph{ample} if a power of it gives an embedding of $X$ in projective space
(Chapter~\ref{kodaira}).
\end{definition}

\begin{example}\label{holomorphic-bundles:example-fubini-study}
On $\CC\PP^n$, the metric on $\mathcal{O}(-1)$ induced from the standard Hermitian metric of $\CC^{n+1}$ has curvature $-\omega_{FS}$ up to a factor, where $\omega_{FS}$ is the Fubini--Study form; hence
$\mathcal{O}(1)$ is positive, with $\frac{i}{2\pi}F = \omega_{FS}$ and $c_1(\mathcal{O}(1)) = [\omega_{FS}]$, the positive generator of $H^2(\CC\PP^n; \ZZ)$
(Example~\ref{characteristic-classes:example-tautological}).
\end{example}

\section{The canonical bundle and adjunction}\label{holomorphic-bundles:section-adjunction}

\begin{proposition}[Adjunction]\label{holomorphic-bundles:proposition-adjunction}
Let $Z \subseteq X$ be a smooth hypersurface with associated line bundle $\mathcal{O}(Z)$. Then the normal bundle is $N_{Z/X} \cong \mathcal{O}(Z)|_Z$, and
\[
K_Z \cong \left(K_X \otimes \mathcal{O}(Z)\right)\big|_Z .
\]
For a smooth hypersurface of degree $d$ in $\CC\PP^n$ one gets $K_Z \cong \mathcal{O}(d - n - 1)|_Z$.
\end{proposition}

\begin{proof}
A local defining function $f_i$ of $Z$ gives, by $df_i$, an isomorphism of the conormal bundle $N^\vee_{Z/X}$ with $\mathcal{O}(-Z)|_Z$: the transition functions of the conormal bundle are the
$f_i/f_j$ restricted to $Z$, computed from $df_i = (f_i/f_j)\,df_j$ modulo $T^*Z$. Taking determinants in the exact sequence
$0 \to N^\vee_{Z/X} \to \Omega^1_X|_Z \to \Omega^1_Z \to 0$ gives $K_X|_Z \cong K_Z \otimes N^\vee_{Z/X} \cong K_Z \otimes \mathcal{O}(-Z)|_Z$, which is the formula. For $\CC\PP^n$ use
$K_{\CC\PP^n} \cong \mathcal{O}(-n-1)$, which follows from the Euler sequence (Example~\ref{vector-bundles:example-euler-sequence}) by taking determinants, and $\mathcal{O}(Z) = \mathcal{O}(d)$.
\end{proof}

\begin{example}\label{holomorphic-bundles:example-adjunction}
A smooth plane curve of degree $d$ in $\CC\PP^2$ has $K_Z \cong \mathcal{O}(d - 3)|_Z$, whence genus $g = (d-1)(d-2)/2$ by Riemann--Roch (Chapter~\ref{curves}); a smooth quadric surface in $\CC\PP^3$ has
$K \cong \mathcal{O}(-2)$ and is $\CC\PP^1 \times \CC\PP^1$; a smooth quartic surface in $\CC\PP^3$ has trivial canonical bundle and is a K3 surface, one of the basic examples in Hodge theory.
\end{example}

\section{Exercises}\label{holomorphic-bundles:section-exercises}

\begin{exercise}\label{holomorphic-bundles:exercise-sections-ok}
Show that the holomorphic sections of $\mathcal{O}(k)$ on $\CC\PP^n$ are the homogeneous polynomials of degree $k$ in the homogeneous coordinates, and that there are none for $k < 0$.
\end{exercise}

\begin{solution}
A section is given by functions $s_i$ on $U_i$ with $s_i = (x_i/x_j)^{-k}s_j$; the function $F(x) = x_i^ks_i(x/x_i)$ is independent of $i$, holomorphic on $\CC^{n+1} \setminus \{0\}$ and homogeneous of degree
$k$. By Hartogs' theorem (Theorem~\ref{several-complex-variables:theorem-hartogs}) it extends to $\CC^{n+1}$, and a homogeneous holomorphic function of degree $k$ is a polynomial for $k \geq 0$ and zero for
$k < 0$, by looking at its power series.
\end{solution}

\begin{exercise}\label{holomorphic-bundles:exercise-degree-line-bundle}
Let $X$ be a compact Riemann surface. Show that $\deg L = \int_Xc_1(L)$ defines a homomorphism $\Pic(X) \to \ZZ$, that $\deg\mathcal{O}(D) = \sum_pa_p$ for $D = \sum a_pp$, and that a line bundle of negative
degree has no nonzero holomorphic sections.
\end{exercise}

\begin{exercise}\label{holomorphic-bundles:exercise-chern-connection-trivial}
Compute the Chern connection and its curvature for the trivial line bundle with metric $e^{-|z|^2}$ on $\CC$, and verify that the curvature is a positive $(1,1)$-form.
\end{exercise}

\begin{exercise}\label{holomorphic-bundles:exercise-picard-torus}
For a complex torus $X = \CC/\Lambda$ of dimension one, show that $H^1(X, \mathcal{O}_X) \cong \CC$ and deduce that $\Pic^0(X)$ is a complex torus of dimension one.
\end{exercise}
